\section{Related Work}\label{sec:related}

\subsection{Collision-Safe VLA Methods}\label{sec:rw_collision}

Recent methods for collision avoidance in VLAs fall into three categories. \textbf{Post-hoc correction} approaches modify policy outputs before execution. AEGIS~\cite{hu2025vlsa} applies Control Barrier Function corrections to VLA actions based on open-set detection and monocular depth estimation, and reports 77.85\% CAR on SafeLIBERO without additional training. iKinQP-RL~\cite{spalter2024ikinqp} enforces collision constraints through inverse-kinematics quadratic programming at runtime. CoFreeVLA~\cite{cofreeevla2026}trains a risk estimator that halts dual-arm execution when predicted collision probability exceeds a threshold.

\textbf{Constrained optimization} embeds safety into training. SafeVLA~\cite{zhang2025safevla}formulates collision avoidance as a Constrained MDP and reports an 83.58\% reduction in safety violations through reward shaping.

\textbf{Inference-time guidance} steers the generative process toward safe trajectories without retraining. Li et al.~\cite{li2024langguided} train an object-centric diffusion policy on 3D point clouds and inject collision costs during denoising. CAPE~\cite{yang2025cape} expands the trajectory distribution through context-aware priors and iterative refinement.

These methods differ in mechanism but share a common evaluation pattern: all report Collision Avoidance Rate as primary evidence of safety. None examine whether the measured CAR reflects genuine obstacle avoidance or conflates it with other behaviors. Our analysis (Section~\ref{sec:fake_safety}) reveals that this conflation is severe in the Pi0 family: 98--99\% of safe episodes are attributable to inaction rather than planned avoidance. We do not challenge the technical contributions of these methods; we identify a validity gap in the metric on which they rely, and propose a filtered variant (CAR$_\text{f}$) that separates genuine avoidance from frozen behavior.

\subsection{Signed Distance Fields in Robot Manipulation}\label{sec:rw_sdf}

Signed Distance Fields encode the distance from a query point to the nearest obstacle surface, where the sign indicates free space (positive) or penetration (negative). DeepSDF~\cite{park2019deepsdf} established neural SDF prediction as a compact shape representation. iSDF~\cite{ortiz2022isdf} extended this to incremental scene reconstruction for real-time robot perception. DiffCo~\cite{zhi2021diffco} introduced differentiable collision checking through learned distance fields and showed that gradient-based trajectory optimization can enforce collision constraints. VoxPoser~\cite{huang2023voxposer} composes 3D value maps from language model outputs for manipulation planning, with avoidance maps that encode obstacle information in volumetric form.

SDF-based collision detection is well validated in classical motion planning. The natural question is whether this effectiveness transfers to VLA action refinement, where the SDF gradient modifies predicted end-effector actions to increase obstacle clearance. Our oracle experiments (Section~\ref{sec:ee_only}) test this directly by replacing the learned SDF network with ground-truth distance fields computed from simulation geometry, which eliminates all prediction error. The results reveal a structural limitation of end-effector-only refinement: all three oracle variants produce more collisions than unrefined baselines, because the gradient that repels the end-effector simultaneously rotates arm links into obstacle volumes. The problem is not SDF accuracy but application scope. Full-body collision avoidance requires distance information along the entire kinematic chain, not only at the tool tip.

\subsection{Evaluation Reliability in Robot Learning}\label{sec:rw_eval}

Concerns about whether standard metrics measure what they claim to measure have a long history in ML evaluation. Benchmark saturation and annotation artifact exploitation in NLP demonstrated that high scores can reflect dataset properties rather than target capabilities~\cite{gururangan2018annotation}. Goodhart's law (a measure ceases to be informative once it becomes an optimization target) has been documented in reinforcement learning settings where agents exploit reward proxies rather than solve intended tasks~\cite{karwowski2024goodhart}. These cases share a common structure: a metric that appears to measure capability $X$ in fact measures a confound $Y$, and the discrepancy remains invisible until the aggregate score is decomposed.

The confound in VLA safety evaluation is action collapse, a phenomenon related to mode collapse in generative models. Diffusion-based robot policies can collapse to the mean of multi-modal action distributions~\cite{chi2023diffusionpolicy},
and VLAs trained on limited demonstration data exhibit similar failure modes that produce near-zero actions. Action collapse itself is well recognized. Its interaction with displacement-based safety metrics is not: a collapsed model registers as maximally safe because it generates no motion that could cause collisions. CAR cannot distinguish a model that plans a collision-free trajectory from a model that does nothing.

All collision avoidance methods reviewed in Section~\ref{sec:rw_collision} adopt CAR without examining its construct validity. Our displacement-filtered variant, CAR$_\text{f}$ (Section~\ref{sec:metrics}), conditions safety measurement on sufficient motor activity and converts CAR from a contaminated aggregate into a metric that isolates genuine avoidance from inaction.
